\documentclass[12pt,a4paper]{article}

\usepackage{amsmath,amssymb,amsthm}
\usepackage{geometry}
\usepackage{bm}
\usepackage{physics}

\geometry{margin=1in}

\begin{document}

\begin{center}
    {\Large \textbf{Quantum Field Theory I}}\\[0.3cm]
    {\large \textbf{First Homework Assignment}}\\[0.5cm]
    \textbf{Name:} Mohammad Parsa Akrami\\
    \textbf{Student Number:} 404210672
\end{center}

\vspace{0.5cm}

\section*{Problem}

We consider a four-dimensional hypercubic lattice with lattice spacing $\Delta$.
The lattice propagator is written as a sum over all lattice paths from $0$ to $x$,
where each link contributes a factor $e^{-m_0\Delta}$. The directions of each
link are
\[
\nu_k\in
\left\{
(\pm1,0,0,0),
(0,\pm1,0,0),
(0,0,\pm1,0),
(0,0,0,\pm1)
\right\}.
\]

The length of a path containing $n$ links is
\[
L=n\Delta.
\]

Therefore, the path sum can be written as
\begin{equation}
A(x)
=
\sum_{n=0}^{\infty}
e^{-m_0\Delta n}
\sum_{\{\nu_k\}}
\delta^{(4)}_{\sum_{k=1}^{n}\nu_k,\,x/\Delta}.
\label{eq:pathsum}
\end{equation}

Here the arguments of the Kronecker delta are dimensionless, since
$x/\Delta$ is the dimensionless lattice coordinate.

\section*{Fourier Representation of the Kronecker Delta}

For dimensionless integer-valued four-vectors, the Kronecker delta can be
represented as
\begin{equation}
\delta^{(4)}_{\bm N,\bm M}
=
\int_{-\pi}^{\pi}
\frac{d^4 k}{(2\pi)^4}
e^{i k\cdot(\bm N-\bm M)}.
\label{eq:kronecker}
\end{equation}

Applying this to Eq.~\eqref{eq:pathsum}, with
\[
\bm N=\sum_{k=1}^{n}\nu_k,
\qquad
\bm M=\frac{x}{\Delta},
\]
gives
\begin{equation}
\delta^{(4)}_{\sum_{k=1}^{n}\nu_k,\,x/\Delta}
=
\int_{-\pi}^{\pi}
\frac{d^4 k}{(2\pi)^4}
\exp\left[
ik\cdot
\left(
\sum_{k=1}^{n}\nu_k-\frac{x}{\Delta}
\right)
\right].
\end{equation}

Thus,
\begin{equation}
A(x)
=
\sum_{n=0}^{\infty}
e^{-m_0\Delta n}
\int_{-\pi}^{\pi}
\frac{d^4 k}{(2\pi)^4}
e^{-ik\cdot x/\Delta}
\prod_{j=1}^{n}
\left(
\sum_{\nu_j}e^{ik\cdot\nu_j}
\right).
\label{eq:fourier-k}
\end{equation}

Since the four-dimensional hypercubic lattice has eight possible directions,
we have
\begin{equation}
\sum_{\nu}e^{ik\cdot\nu}
=
2\left[
\cos k_1+\cos k_2+\cos k_3+\cos k_4
\right].
\label{eq:sumdirections}
\end{equation}

Therefore,
\begin{equation}
A(x)
=
\sum_{n=0}^{\infty}
e^{-m_0\Delta n}
\int_{-\pi}^{\pi}
\frac{d^4 k}{(2\pi)^4}
e^{-ik\cdot x/\Delta}
\left[
2\sum_{\mu=1}^{4}\cos k_\mu
\right]^n.
\label{eq:Ak}
\end{equation}

Now $k$ is a \emph{dimensionless} lattice momentum. To introduce the
physical momentum $p$, we define
\begin{equation}
k_\mu=p_\mu\Delta,
\qquad
p_\mu=\frac{k_\mu}{\Delta}.
\end{equation}

Consequently, the integration range becomes
\[
-\frac{\pi}{\Delta}\leq p_\mu\leq\frac{\pi}{\Delta},
\]
and the exponent becomes
\[
e^{-ik\cdot x/\Delta}=e^{-ip\cdot x}.
\]

Thus, up to the overall normalization factor coming from
$d^4k=\Delta^4d^4p$, we can write
\begin{equation}
A(x)
\sim
\sum_{n=0}^{\infty}
e^{-m_0\Delta n}
\int
d^4p\,
e^{-ip\cdot x}
\left[
2\sum_{\mu=1}^{4}\cos(p_\mu\Delta)
\right]^n.
\label{eq:Ap}
\end{equation}

For small $\Delta$,
\begin{equation}
2\sum_{\mu=1}^{4}\cos(p_\mu\Delta)
=
8-p^2\Delta^2+O(\Delta^4),
\end{equation}
where
\[
p^2=p_1^2+p_2^2+p_3^2+p_4^2.
\]

Hence,
\begin{align}
\left[
2\sum_{\mu=1}^{4}\cos(p_\mu\Delta)
\right]^n
&\simeq
\left(8-p^2\Delta^2\right)^n
\nonumber\\
&=
8^n
\left(
1-\frac{p^2\Delta^2}{8}
\right)^n
\nonumber\\
&\simeq
8^n
e^{-np^2\Delta^2/8}.
\label{eq:gaussian}
\end{align}

Substituting this into Eq.~\eqref{eq:Ap},
\begin{equation}
A(x)
\sim
\sum_{n=0}^{\infty}
e^{-m_0\Delta n}
8^n
\int d^4p\,
e^{-np^2\Delta^2/8}
e^{-ip\cdot x}.
\end{equation}

Therefore,
\begin{equation}
A(x)
\sim
\sum_{n=0}^{\infty}
\exp\left[
-n\left(m_0\Delta-\ln 8\right)
\right]
\int d^4p\,
e^{-np^2\Delta^2/8}
e^{-ip\cdot x}.
\label{eq:An}
\end{equation}

We see that the important combination of $\Delta$ and $n$ is
$n\Delta^2$, rather than $n\Delta$. Thus, in the continuum limit, we define
\begin{equation}
s=\frac{n\Delta^2}{8}.
\label{eq:sdef}
\end{equation}

In order for long paths to remain relevant in the continuum limit, the
factor multiplying $n$ in Eq.~\eqref{eq:An} must be of order $\Delta^2$.
Therefore, we tune the bare mass according to
\begin{equation}
m_0\Delta-\ln 8
=
\frac{m^2\Delta^2}{8}.
\label{eq:mass-tuning}
\end{equation}

This gives
\begin{equation}
\boxed{
m_0(\Delta)
=
\frac{\ln 8}{\Delta}
+
\frac{m^2}{8}\Delta
}.
\label{eq:m0}
\end{equation}

With this choice,
\begin{equation}
e^{-n(m_0\Delta-\ln8)}
=
e^{-m^2n\Delta^2/8}
=
e^{-m^2s},
\end{equation}
and therefore
\begin{equation}
A(x)
\sim
\sum_n
e^{-m^2s}
\int d^4p\,
e^{-sp^2}
e^{-ip\cdot x}.
\end{equation}

In the continuum limit, the sum over $n$ becomes an integral over $s$, up to
an overall normalization factor:
\begin{equation}
A(x)
\sim
\int_0^\infty ds\,
e^{-m^2s}
\int d^4p\,
e^{-sp^2}
e^{-ip\cdot x}.
\end{equation}

Thus the required continuum limit is obtained by tuning the bare mass as in
Eq.~\eqref{eq:m0}.
\section*{Shortest Path Domination for Finite $\Delta$}

For finite $\Delta$, there is another regime in which the shortest paths
dominate the path sum.

Each additional link has a factor
\[
e^{-m_0\Delta},
\]
while at every step there are eight possible directions. Thus, very roughly,
the competition between the number of paths and the suppression of each path
is controlled by
\begin{equation}
8e^{-m_0\Delta}.
\end{equation}

If
\begin{equation}
8e^{-m_0\Delta}<1,
\end{equation}
or equivalently
\begin{equation}
m_0\Delta>\ln8,
\label{eq:shortestcondition}
\end{equation}
longer paths are exponentially suppressed.

For a fixed endpoint $x$, the shortest possible path has the minimum number
of links,
\begin{equation}
n_{\min}
=
\sum_{\mu=1}^{4}
\left|\frac{x_\mu}{\Delta}\right|.
\end{equation}

Any longer path must contain additional steps. In particular, because the
lattice is bipartite, paths connecting the same two points differ in length
by an even number of links, so a path with length
\[
n=n_{\min}+2r
\]
contains $2r$ additional links compared with a shortest path.

Each additional pair of links introduces an exponential suppression of the
form
\[
e^{-2m_0\Delta}.
\]

Therefore, when
\[
m_0\Delta\gg\ln8,
\]
the exponential suppression of longer paths overwhelms the growth in the
number of available paths, and the path sum is dominated by the shortest
paths.

Hence, for finite $\Delta$, the two relevant regimes can be summarized as
\begin{equation}
\begin{array}{ccl}
m_0\Delta>\ln8
&\Longrightarrow&
\text{long paths are suppressed,}\\[0.2cm]
m_0\Delta\simeq\ln8
&\Longrightarrow&
\text{long paths can remain relevant.}
\end{array}
\end{equation}

In particular, the continuum limit requires tuning $m_0\Delta$ to approach
$\ln8$ with a difference of order $\Delta^2$, as in
Eq.~\eqref{eq:mass-tuning}.

For a fixed finite $m_0$ while taking $\Delta\to0$, however,
\[
m_0\Delta\longrightarrow0,
\]
and therefore
\[
8e^{-m_0\Delta}\longrightarrow8>1.
\]
Thus the path sum is not in the shortest-path regime in this limit; instead,
the number of long paths grows too rapidly and the sum diverges. The
shortest-path domination discussed above applies to finite $\Delta$ when
$m_0\Delta$ is sufficiently larger than $\ln8$.

\section*{AI use declaration}

The AI model used in preparing and revising this homework solution was
\textbf{GPT-5.6 Luna}.

The solution was checked using the course lecture notes as a reference,
in particular the derivation of the lattice path sum, the Fourier
representation, and the continuum limit involving the combination
$n\Delta^2$.

The AI was also used to revise and correct a small mistake in the
calculation: the Kronecker delta was rewritten with dimensionless
arguments, and the distinction between the dimensionless lattice momentum
$k$ and the physical momentum $p$ was made explicit. The main derivation
and notation of the original solution were otherwise retained.\\
The AI was also used to convert the hand written solution into Latex code.
\end{document}